% MSR 2027 -- Technical Papers
% Format: IEEE conference proceedings (24pt title, 10pt text).
% DO NOT use the compsoc / compsocconf options.
% Limit: 10 pages of main text (including figures, tables, and appendices)
%        + 2 pages containing only references.
% Double-anonymous review: do not include names, affiliations, or acknowledgments
% that identify the authors in the submission version.
\documentclass[10pt,conference]{IEEEtran}

\usepackage{cite}
\usepackage{amsmath,amssymb,amsfonts}
\usepackage{algorithmic}
\usepackage{graphicx}
\usepackage{textcomp}
\usepackage{xcolor}
\usepackage{url}

\begin{document}

\title{The Vocabulary of Flaky Tests in C\#}

% Anonymous submission (double-anonymous). Restore the authors only in the camera-ready version.
\author{\IEEEauthorblockN{Anonymous Author(s)}
\IEEEauthorblockA{\textit{Omitted for double-anonymous review}}
}

\maketitle

\begin{abstract}
Flaky tests produce non-deterministic outcomes without any change to the code, eroding trust in continuous integration. Vocabulary-based prediction has proven effective for Java, JavaScript, and Swift but, to the best of our knowledge, has never been evaluated for C\#, whose concurrency model (\texttt{Task}, \texttt{async}/\texttt{await}) and testing frameworks differ from those previously studied. We re-ran the test suites of 17 open-source C\# projects, totaling 19,211 test cases, 100 times, and a further 20 times in isolation. After grouping by method, we obtained 30 flaky and 12,635 stable test methods. We trained eight classifiers on TF-IDF features of unigrams and bigrams. The best ones, Extra Trees and Random Forest, reach an MCC of 0.51--0.53 and an AUC of 0.86--0.87 in cross-validation, above trivial baselines but only moderately above a vocabulary-threshold rule (MCC 0.40) and a structure-only model (MCC 0.36). Across repositories, AUC drops to 0.75--0.78 and precision to 2--3\%: the approach is suited to prioritizing reruns, not to blocking pipelines. The markers of flakiness are Task-based concurrency, time handling, and resource lifecycles, and we found no robust marker of stability. The model fails when the root cause lies outside the test body or when asynchronous constructs are used deterministically, which exposes the limited reach of vocabulary.
\end{abstract}

\begin{IEEEkeywords}
flaky tests, test flakiness, software testing, vocabulary-based prediction, C\#, .NET
\end{IEEEkeywords}

\section{Introduction} 
In software development, tests are essential for identifying defects and verifying that fixes work correctly \cite{kapfhammer2004software}. Ideally, test results remain stable across runs of an unchanged codebase. However, flaky tests, that is, tests that pass or fail nondeterministically without any change to the code or the execution environment, are a common reality \cite{pinto2020vocabulary}. Their causes are diverse, including concurrency issues, asynchronous operations, order dependencies between tests, improper management of external resources, dependence on the system clock, and the use of random values \cite{luo2014empirical}, as well as overly narrow ranges of expected values and timeout-related issues \cite{eck2019understanding}.
Although flaky tests do not directly affect software functionality, they impose a substantial cost on delivery pipelines. Developers report wasted effort analyzing failures, ignored failures, and loss of trust in the test suite \cite{eck2019understanding, gruber2022survey}, and at Google 16\% of tests exhibited flaky behavior, consuming up to 16\% of testing capacity in reruns \cite{micco2017state}. Detecting flakiness through reruns is expensive, since tests that fail rarely require many executions \cite{lam2020study}. Static prediction from the test code is therefore attractive: it requires running nothing and produces a human-readable explanation.
Pinto et al. \cite{pinto2020vocabulary} showed that flaky tests share a characteristic vocabulary, that is, identifiers and keywords that occur disproportionately in them, and achieved an F-score of 0.95 in Java. The idea was replicated for JavaScript \cite{soratto2023vocabulary}, extended to five languages \cite{ahmad2025understanding}, and applied to Swift \cite{medeiros2026vocabulary}. To the best of our knowledge, it has never been evaluated for C\#, even though the language has relevant particularities. Its concurrency model centers on the Task-based asynchronous pattern (Task, async/await, CancellationToken) and the thread pool, which introduces its own timing and scheduling dependencies. In tests from Microsoft .NET projects, concurrency is the leading cause of flakiness, with fixed waits and timeouts as a frequent pattern \cite{leesatapornwongsa2022flakerepro}. Test suites are written with several testing frameworks (xUnit, NUnit, MSTest) on two execution platforms, and libraries are often compiled for multiple target frameworks, so the same test may run under different runtimes. Integration suites also depend on containerized services and on environment behaviors such as locale, time zone, and file system paths. If the lexical patterns that signal flakiness in Java, JavaScript, and Swift do not transfer to C\#, vocabulary-based tools built around them would offer little help to .NET developers. This is the gap this work addresses.
To address it, we investigate static prediction of flaky tests in C\# from the vocabulary of test source code. Our approach, organized in six steps (Section 4), selects open-source C\# projects and systematically reruns their test suites, 100 times in an automated fashion and again 20 times under controlled isolation; labels each test as flaky, broken, or stable according to its outcomes; resolves inherited tests to the method that declares them; extracts the vocabulary and structural features of each test with the Roslyn compiler platform; and trains classifiers on this vocabulary to predict a test's propensity to be flaky without executing it, while also extracting an interpretable ranking of the tokens most strongly associated with flakiness.
We applied this methodology to 17 projects, with 19,211 test cases and 12,665 labeled test methods, of which 30 are flaky and 12,635 are stable (0.24\%). Among the eight classifiers evaluated, tree ensembles performed best under cross-validation, with Extra Trees reaching an MCC of 0.53 and an AUC of 0.86, and Random Forest an MCC of 0.51 and an AUC of 0.87. These results are well above random guessing (MCC 0.02), but only moderately above a structure-only baseline (MCC 0.36) and a vocabulary threshold rule (MCC 0.40). When each repository is held out and evaluated at its natural base rate (0.39\%), AUC drops to 0.75--0.78 and precision stays at 2--3\%; thus, the approach is better suited to ranking tests and prioritizing reruns than to gating a pipeline: the top quartile of each repository's ranking contains 23 of the 30 flaky tests.
The main contributions of this work are:
A dataset of flaky, broken, and stable C\# tests, built from 100 automated and 20 controlled reruns of 17 open-source projects, with inherited tests resolved to the methods that declare them.
An evaluation of eight classifiers for vocabulary-based static prediction in C\#, under repeated cross-validation and under cross-repository validation at the natural base rate, showing that the approach transfers to C\#, with clear limits.
A C\# flakiness vocabulary, which separates frequent general markers from rare, specific ones and shows that no robust marker of stability emerges.
An analysis of the model's errors in the test code, with sensitivity analyses without attributes and with helper-method expansion, which outline the limits of lexical prediction.
An assessment of label robustness, comparing the two rerun rounds (which agree on only 13 of the 36 flaky tests), and documentation of the pitfalls of rerun studies in .NET that distorted our initial counts.

\section{Motivating Example}

Consider ReconnectOnRouterBug, a test from NetMQ, a messaging library in the .NET ecosystem, which was flaky in our 100-run round. The test checks a reconnection behavior. It connects a router socket to a dealer, sends a message, disconnects, reconnects, and sends the message again. The goal is to ensure that the second message also reaches its destination after reconnection (abbreviated in Figure 1).

\begin{figure}[!htbp]
\noindent\rule{\columnwidth}{0.4pt}\\[3pt]
{\footnotesize
\begin{tabular}{@{}r@{\hspace{8pt}}l@{}}
{\scriptsize 1} & \verb|router.Connect(addr);    Thread.Sleep(100);| \\
{\scriptsize 2} & \verb|router.SendMoreFrame("dealer")| \\
{\scriptsize 3} & \verb|      .SendFrame("Hello");| \\
{\scriptsize 4} & \verb|router.Disconnect(addr); Thread.Sleep(1000);| \\
{\scriptsize 5} & \verb|router.Connect(addr);    Thread.Sleep(100);| \\
{\scriptsize 6} & \verb|router.SendMoreFrame("dealer")| \\
{\scriptsize 7} & \verb|      .SendFrame("Hello");| \\
\end{tabular}}\\[1pt]
\rule{\columnwidth}{0.4pt}
\caption{Flaky test with fixed waits (\texttt{ReconnectOnRouterBug}, abbreviated).}
\label{fig:flaky}
\end{figure}

The key point is that the test has no mechanism for synchronizing with what happens inside the library. Connecting and disconnecting a socket are asynchronous operations: the Connect call returns before the connection is actually established. The test works around this with fixed waits (Thread.Sleep(100) after each connection and Thread.Sleep(1000) after disconnection). Each of these waits is a bet on how long the connection takes. If the bet holds in most runs, the test passes. But in a slower or more heavily loaded run, such as on a shared continuous integration machine, the connection may take longer than the wait. The second send then reaches the socket before the reconnection is reestablished, and the test fails, without anything in the code having changed.

This kind of failure is hard to handle with dynamic tools. It depends on a sub-second time window and on load conditions we do not control, and therefore it does not reproduce on demand. Confirming that the test is flaky required repeating the suite 100 times, and a smaller sample might not have observed the failure. The situation is that of a developer who sees the test fail once in continuous integration, reruns it, sees it pass, and dismisses the failure as noise. This cycle is the wasted effort and loss of trust described in Section 1.

Even so, the likely cause is visible in the source code. The test is dominated by socket calls (Connect, Disconnect, SendFrame) interleaved with Thread.Sleep, and it contains no wait for an event, such as confirmation that the connection has been established. This is the starting point of our work: a failure that manifests only at run time can leave a lexical trace in the test's text. If this trace exists and is consistent, a model trained only on the code's terms can point to the most suspicious tests without running them. Section 6.5 returns to this test for a complementary reason: the model ranked it low. The trace here is a synchronous Sleep, not the async/await vocabulary that dominates the flaky tests in our dataset, which suggests that the model recognizes one kind of cause better than another.

\section{Related Work}
Flaky tests. A flaky test produces different outcomes across runs of the same software version \cite{luo2014empirical}. Developers spend considerable effort triaging failures that do not signal a real defect \cite{pinto2020vocabulary}, and repeated exposure erodes trust to the point that genuine failures are ignored \cite{gruber2022survey, rahman2018impact}; the cost to delivery pipelines is also documented \cite{herzig2015empirically, micco2017state}.
Detection. Techniques are dynamic, hybrid, or static, depending on whether they require running the test. Rerunning the suite is practiced at scale at Google and Microsoft \cite{lam2020study, micco2017state}, but its cost grows with the rerun budget, and even five reruns reveal only about 88\% of flaky tests \cite{lam2020understanding}. DeFlaker flags tests that fail without exercising modified code \cite{bell2018deflaker}; Shaker injects stress to expose concurrency-related flakiness \cite{silva2020shake}. For .NET, FlakeRepro \cite{leesatapornwongsa2022flakerepro} combines static and dynamic analysis to reproduce the failure of a concurrency-related flaky test: across 31 flaky tests from 22 Microsoft projects, it reproduced 26. It requires the test binaries and an observed failure, so it acts during debugging after flakiness has been found, whereas our approach predicts which tests are likely flaky from source code alone and can feed into it.
Vocabulary-based prediction. Pinto et al.\ \cite{pinto2020vocabulary} trained classifiers on the vocabulary of test code (F1 = 0.95 in Java). Subsequent work confirmed its robustness across feature sets and classifiers (FlakeFlagger \cite{alshammari2021flakeflagger}, FLAST \cite{verdecchia2021know}) and across languages (JavaScript \cite{soratto2023vocabulary}, five languages \cite{ahmad2025understanding}); Barbosa et al.\ \cite{barbosa2022test} show that root causes vary by ecosystem. Medeiros et al.\ \cite{medeiros2026vocabulary} studied Swift and found a vocabulary of asynchronous constructs and stability markers. No study covers C\#.
Positioning. This work differs from previous ones in four respects. First, the unit of analysis is the declaring method, with inheritance resolved, not the test case. Second, we compare two rerun rounds and measure how much the labels depend on them. Third, in addition to balanced cross-validation, we evaluate across repositories at the natural rate, with average precision and a ranking analysis. Fourth, we separate frequent, general markers from rare, specific ones, instead of treating the vocabulary as a single ranking.
Learned features. Pre-trained models such as CodeBERT \cite{fatima2022flakify} and prompting-based use of LLMs \cite{berndt2026can} provide mixed evidence: learned features help more in categorizing root causes than in binary prediction \cite{rahman2024quantizing}, and accuracy drops considerably under realistic imbalance \cite{rahman2025understanding}. Vocabulary features remain cheap and interpretable.

\section{Approach}
Our approach rests on one premise: although flakiness manifests only at run time, it leaves a stable lexical trace in the test code. Figure 2 summarizes the six-step pipeline.

\begin{figure*}[!htbp]
\centering
{\scriptsize
\setlength{\fboxsep}{3pt}%
\newcommand{\etapa}[2]{\fbox{\parbox[c][11ex][c]{\dimexpr(\textwidth-55pt)/6-2\fboxsep-2\fboxrule\relax}{\centering\hyphenpenalty=10000\textbf{#1}\\[2pt]#2}}}%
\etapa{1. Subjects}{17 open-source\\C\# projects}%
\hfill$\rightarrow$\hfill
\etapa{2. Reruns}{100$\times$ + 20$\times$\\isolated;\\TRX reports}%
\hfill$\rightarrow$\hfill
\etapa{3. Labels}{flaky, broken,\\stable, skipped;\\declaring method}%
\hfill$\rightarrow$\hfill
\etapa{4. Processing}{Roslyn: tokens\\and structure}%
\hfill$\rightarrow$\hfill
\etapa{5. Models}{TF-IDF;\\8 classifiers;\\1:1; CV and LOPO}%
\hfill$\rightarrow$\hfill
\etapa{6. Vocabulary}{information\\gain of\\tokens}%
}
\caption{End-to-end pipeline for building the vocabulary of flaky tests in C\#.}
\label{fig:metodologia}
\end{figure*}

Step 1: Project selection. We started from open-source C\# libraries hosted on GitHub with extensive test suites and then filtered them by an objective criterion: the project had to build and run its entire suite in our environment. The result is 17 projects and 19,211 test cases, covering messaging, actor systems, scheduling, RPC, logging, resilience, and architecture testing. Two candidate projects (StackExchange.Redis and SSH.NET) were excluded because their suites depend on external servers whose runs were disrupted by the environment.
Step 2: Reruns. Each suite was built once and run 100 times in a row, without rebuilding between runs. In each run, the .NET test runner (dotnet test) produced a TRX (Test Results XML) report, an XML file recording each executed test. We built with dotnet build and ran with dotnet test \texttt{--no-build} and the TRX logger; xUnit v3 projects on Microsoft.Testing.Platform were run through the test executable itself, with \texttt{--report-trx}. These reports were the basis for identifying tests whose outcomes changed across runs. Some rounds suffered environmental disruptions, such as leftover containers from previous runs and a degraded Redis cluster. They were discarded and repeated in isolation. Since flakiness is rare (fewer than 1\% of tests), a few unstable tests are enough to change the results, and any environmental noise could be mistaken for flakiness. We therefore performed a second round of 20 runs under controlled isolation (5 for Testcontainers, whose suite requires containers) to verify the tests flagged in the first. Everything ran on a laptop with an Intel Core i7-1165G7, 16 GB of RAM, and Windows 11. The 100-run rounds took place in July 2026, before a .NET SDK update whose previous version we did not record; the 20-run rounds used the SDK resolved by each repository's global.json. This is a limitation of the 100-run labels. The SDKs installed side by side are 8.0.425, 9.0.318, 10.0.112, 10.0.303, and 10.0.401; the 20-run round used 10.0.401 in 11 repositories, 10.0.303 in Nerdbank.Streams, Polly, and StreamJsonRpc, 10.0.112 in Akka.NET, and 8.0.425 in the RabbitMQ client.
Step 3: Labeling. Each test received a label according to its outcomes across runs: flaky (passes and fails), broken (always fails), stable (never fails), or skipped (skipped in every run). Broken and skipped tests were excluded from modeling, which compares only flaky against stable tests. The unit of analysis is the declaring test method, not the test case reported by the runner. Parameterized tests and methods inherited from base classes generate several cases from a single method. Without this grouping, the same code would be counted several times, and nearly identical tests would fall into different cross-validation folds. The choice of unit changes the count: the 30 flaky methods correspond to 50 flaky test cases, because some methods generate several cases (through parameterization or inheritance). Since test identity in TRX files is fragile, we identified each test by its fully qualified class plus its display name and handled the pitfalls in Table 1: we normalized timestamps, thread identifiers, and run dates in the names, ran both rounds on the same target framework, and discarded contaminated rounds.

\begin{table}[!htbp]
\caption{Methodological pitfalls of rerun studies in .NET.}
\label{tab:armadilhas}
\centering
\footnotesize
\renewcommand{\arraystretch}{1.15}
\begin{tabular}{@{}p{0.30\columnwidth}p{\dimexpr0.70\columnwidth-2\tabcolsep\relax}@{}}
\hline\hline
\raggedright \textbf{Pitfall} & \raggedright \textbf{Example and effect} \tabularnewline
\hline
\raggedright Nondeterministic names & \raggedright Akka.NET \emph{theories} embed timestamps and \emph{thread} ids: 3,308 distinct names for 1,328 tests \tabularnewline
\raggedright Dates in names & \raggedright NLog \emph{theories} embed the current date: 2,876 vs.\ 2,870 tests across rounds \tabularnewline
\raggedright Target framework & \raggedright Nerdbank.Streams: 1,142 tests on \texttt{net472} vs.\ 877 on \texttt{net8.0} \tabularnewline
\raggedright Name collisions & \raggedright In some projects, TRX stores only the method name: tests from different classes get mixed \tabularnewline
\raggedright Contaminated environment & \raggedright Leftover containers and a degraded \emph{cluster} produce spurious flaky tests \tabularnewline
\raggedright Inherited tests & \raggedright 363 of 934 methods (Nerdbank.Streams) and 2,199 of 2,495 (StreamJsonRpc) are inherited \tabularnewline
\hline\hline
\end{tabular}
\end{table}

Step 4: Test processing. Using Roslyn, the C\# compiler, we extracted each test method from the source code, identifying it by its attributes ([Fact], [Test], [TestMethod], and equivalents). From each method, we recorded the tokens (identifiers and keywords), discarding reserved words and common language modifiers (such as public, static, var, if, and return), which appear in almost every test and do not distinguish them. We kept async and await, which indicate asynchronous code. We also recorded size and structural features: number of lines, number of awaits and assertions, and whether the method is async. Since many tests are inherited, we resolved each executed test to its declaring method, which is flaky if any inherited instance is flaky.

Step 5: Model construction. The tokens form the test's vocabulary and are vectorized with TF-IDF (term frequency--inverse document frequency), a technique that turns the text of each test into a numeric vector. Each term receives a weight that combines how often it appears in that test (TF) and the inverse of the proportion of tests that contain it (IDF). Thus, the weight is high when the term appears often in the test and is rare in the set (such as weakreference) and low when it is present in almost every test (such as assert). We used single terms (unigrams, such as await) and pairs of adjacent terms (bigrams, such as await task), with sublinear term frequency, which reduces the weight of heavily repeated terms. We kept only terms present in at least two tests and in at most 90\% of them. The vectorizer is fitted only on the training tests of each fold, which prevents information leakage from the evaluation tests. We trained eight classifiers: Random Forest and Extra Trees (200 trees each), Gradient Boosting, Decision Tree, Naive Bayes, linear SVM, KNN, and logistic regression. The first three are ensemble models, which combine many decision trees. To find out what the vocabulary actually adds, we compared all of them against four baselines: a majority predictor, a random predictor, a threshold rule based on the 10 tokens with the highest information gain, and a Random Forest trained only on the size and structural features, without any tokens.
Step 6: Vocabulary extraction. Finally, we ranked the tokens by information gain with respect to the label, which measures how much knowing whether the token is present reduces uncertainty about the test's class, in nats and averaged over 30 subsampling draws, together with its frequency in each class, the number of repositories in which it occurs in flaky tests, and the lift between classes, that is, how many times more frequent the token is among flaky tests than among stable ones. This ranking makes the approach explanatory rather than opaque: it allows the predictive signal to be linked to concrete constructs and to the known root causes of flakiness.

\section{Objects of Analysis}
Table 2 lists the 17 projects, which span serialization, logging, resilience, messaging, actor, scheduling, RPC, and architecture testing libraries. They are libraries distributed as NuGet packages, so their tests exercise public APIs rather than user interfaces or deployed services. Most are oriented toward concurrency, asynchrony, or I/O, the domains in which timing and ordering dependencies are most likely; Serilog, Newtonsoft.Json, and ArchUnitNET are the main exceptions and serve as contrast projects. Beyond C\#, the projects share other characteristics. They are mature, widely used projects with long histories and large suites (from 76 tests in MassTransit.Abstractions to 3,456 in Newtonsoft.Json), whose tests live in the same repository and run with the .NET testing tools. Several target multiple frameworks (for example, net472 and net8.0 in Nerdbank.Streams), which affects test identity (Table 1). All passed the same executability filter in Step 1. As for testing frameworks, xUnit is the most common, used in 13 of the 16 repositories; MassTransit, Quartz.NET, and Newtonsoft.Json use NUnit. Because the set was assembled by convenience, it does not represent the population of .NET projects, and this limitation is discussed in Section 8.

\begin{table}[!h]
\caption{Studied projects. Counts of flaky tests per test case and per declaring method (units) in the 100-run and 20-run rounds. Brk.: broken tests.}
\label{tab:projetos}
\centering
\scriptsize
\setlength{\tabcolsep}{2.7pt}
\begin{tabular}{@{}lrrrrrr@{}}
\hline\hline
 & & \multicolumn{2}{c}{\textbf{Flaky 100$\times$}} & \multicolumn{2}{c}{\textbf{Flaky 20$\times$}} & \textbf{Brk.} \\
\textbf{Project} & \textbf{Tests} & \textbf{cases} & \textbf{units} & \textbf{cases} & \textbf{units} & \textbf{100$\times$} \\
\hline
Serilog                  &    808 &  0 & 0 &  0 & 0 &  0 \\
Testcontainers-dotnet    &    507 &  1 & 1 &  0 & 0 &  0 \\
Polly                    &    774 &  3 & 3 &  3 & 3 & 14 \\
Newtonsoft.Json          &  3,456 &  0 & 0 &  0 & 0 &  4 \\
NLog                     &  2,870 &  2 & 2 &  5 & 5 &  3 \\
MassTransit.Abstractions &     76 &  1 & 1 &  1 & 1 &  0 \\
MassTransit.Tests        &  1,133 &  7 & 5 &  7 & 4 &  6 \\
RestSharp                &    564 &  0 & 0 &  0 & 0 &  0 \\
AsyncEx                  &    266 &  0 & 0 &  0 & 0 &  0 \\
Akka.NET                 &  1,328 &  1 & 1 &  1 & 1 &  8 \\
RabbitMQ (unit)          &    127 &  0 & 0 &  0 & 0 &  0 \\
Coravel                  &  1,008 &  0 & 0 &  0 & 0 &  0 \\
Quartz.NET               &  1,388 &  2 & 1 &  1 & 0 & 10 \\
Nerdbank.Streams         &  1,142 &  6 & 6 &  0 & 0 &  3 \\
StreamJsonRpc            &  2,837 & 25 & 8 & 11 & 5 &  5 \\
NetMQ                    &    286 &  2 & 2 &  0 & 0 &  0 \\
ArchUnitNET              &    641 &  0 & 0 &  0 & 0 &  2 \\
\hline
\textbf{Total}           & \textbf{19,211} & \textbf{50} & \textbf{30} & \textbf{29} & \textbf{19} & \textbf{55} \\
\hline\hline
\end{tabular}
\end{table}


Broken tests, which fail in every run and are excluded from modeling, total 55 cases in the 100-run round, concentrated in Polly (14) and Quartz.NET (10). The 19,211 test cases correspond to 15,705 distinct test methods, counting inherited instances, and to 12,797 declaring methods after resolving inheritance; of these, 12,665 are flaky or stable and form the modeling set: 30 flaky and 12,635 stable (0.24\%) with the 100-run labels, and 19 flaky and 12,645 stable with the 20-run labels. Two properties shape the evaluation. First, the set is severely imbalanced. Second, flakiness is concentrated: only nine repositories have any flaky test, and three of them (StreamJsonRpc, with 8 units, MassTransit, with 6 across two suites, and Nerdbank.Streams, with 6) account for 20 of the 30. StreamJsonRpc's 25 flaky cases come from only 8 declaring methods because the same test is run with six serializers.

\section{Evaluation}
We structure the evaluation around four research questions, the first with two subquestions.
RQ1.1 --- Predictive effectiveness. How accurately can flaky tests in C\# be predicted from the vocabulary of their source code?
RQ1.2 --- Label robustness. Do the labels and results remain stable between the 100-run and the 20-run rounds?
RQ2 --- Baselines. Does vocabulary-based prediction outperform the majority, random, vocabulary-threshold, and structure-only baselines?
RQ3 --- The vocabulary. Which tokens are most strongly associated with flakiness?
RQ4 --- Error analysis. When and why does the classifier fail?
Rationale. RQ1.1 establishes whether the lexical trace of flakiness is strong enough, in C\#, for standard classifiers to separate flaky from stable tests using only static features, since the method has never been evaluated for a language with this concurrency model. RQ1.2 tests the labeling strategy: the conclusions of RQ1.1 hold only if the labels, which come from reruns and therefore from a nondeterministic process, are reasonably stable. RQ2 puts the results in context against trivial and structural baselines, to ensure that performance reflects genuine discriminative power rather than an artifact of the evaluation. RQ3 addresses interpretability: which lexical units drive the prediction, so they can be related to known root causes. RQ4 examines representative errors, to understand the limits of lexical prediction and when a static detector should or should not be trusted.

\subsection{Experimental setup}
Because the dataset is highly imbalanced (0.24\% flaky), a model that answered ``stable'' for everything would be right more than 99\% of the time without learning anything. We therefore evaluate in two stages. In the first, we train on a random undersample of the stable tests, so that each round has as many stable tests as flaky ones (1:1 ratio), and apply stratified 5-fold cross-validation, which uses, each time, four groups of tests for training and one for evaluation, keeping the same proportion of flaky tests in all of them. We repeat the process with 10 different draws, yielding 50 estimates per model, and report the mean with a 95\% interval, which reflects variation across draws and folds but not uncertainty about other projects. Since the evaluation folds are also balanced, precision and F1 are inflated (a predictor that always answers ``flaky'' already obtains an F1 of 0.67). We therefore compare models by MCC, the Matthews correlation coefficient, which summarizes hits and errors in both classes in a single number from $-1$ to 1 (0 corresponds to chance and 1 to getting everything right), and by AUC, the area under the ROC curve, which shows how the model trades off flaky tests detected against stable tests wrongly flagged at each decision threshold. AUC equals the probability that the model gives a higher score to a randomly drawn flaky test than to a randomly drawn stable one: it is 0.5 when the model ranks like chance and 1.0 when it ranks perfectly. Because it evaluates the ordering without depending on a threshold, it suits the use we propose, which is ranking tests to decide which to rerun first.
In the second stage, we evaluate what happens when facing a new project, since cross-validation mixes tests from all repositories and may reward the model merely for recognizing conventions of a project it has already seen. We use the leave-one-repository-out (LOPO) scheme: we train on all repositories but one and test on all tests of the held-out repository, at the natural proportion of 0.39\% flaky tests. In this scenario, the closest to real use, we report average precision (AP) in addition to AUC. AP walks down the ranking from top to bottom and computes, at each flaky test found, what fraction of the tests seen so far is actually flaky; the result is the mean of these fractions. It is high when flaky tests are concentrated at the top of the list and low when they are mixed among many stable ones. Unlike AUC, AP is sensitive to the rarity of positives: a model with no ranking power obtains an AP equal to the proportion of flaky tests, about 0.004. The difference between the results of the two stages indicates how much of the performance comes from general flakiness patterns and how much comes from each project's particular vocabulary.
As for implementation, Random Forest and Extra Trees use 200 trees; the SVM uses a linear kernel; KNN uses $k = 5$ with cosine distance; the remaining classifiers use scikit-learn's default values. In LOPO, training uses the same 1:1 undersampling of the other repositories (10 draws), the scores are the mean over the 10 draws, and the held-out repository is evaluated in full, at its natural class distribution. We report five complementary metrics in cross-validation (precision, recall, F1, MCC, and AUC) and four in LOPO (AUC, AP, recall, and precision).

\subsection{RQ1.1 and RQ2: predictive effectiveness and baselines}
In cross-validation, the ensemble models performed best: Extra Trees, Random Forest, and Gradient Boosting reach an MCC of 0.51--0.53 and an AUC of 0.83--0.87 (Table 3). The other classifiers range from 0.38 to 0.44 MCC, and the random and majority baselines stay at chance level. Two non-trivial baselines, however, are strong: the rule that uses only the ten most informative tokens reaches an MCC of 0.40, and the Random Forest that uses only the test's structure (without vocabulary) reaches 0.36. Vocabulary therefore adds about 0.15--0.17 MCC for the best models, and less for the others (from 0.02 to 0.08). Part of the threshold rule's strength comes from project-specific tokens (Section 6.4), which inflates it under cross-validation but not across repositories.

\begin{table}[!h]
\caption{Cross-validation with 100-run labels (mean $\pm$ 95\% CI; balanced test folds).}
\label{tab:cv100}
\centering
\scriptsize
\setlength{\tabcolsep}{2.5pt}
\begin{tabular}{@{}p{46pt}ccccc@{}}
\hline\hline
\textbf{Model} & \textbf{Prec.} & \textbf{Rec.} & \textbf{F1} & \textbf{MCC} & \textbf{AUC} \tabularnewline
\hline
\raggedright Extra Trees            & 0.87${\pm}$0.04 & 0.63${\pm}$0.06 & 0.69${\pm}$0.04 & 0.53${\pm}$0.05 & 0.86${\pm}$0.03 \tabularnewline
\raggedright Gradient Boosting      & 0.80${\pm}$0.04 & 0.70${\pm}$0.05 & 0.72${\pm}$0.04 & 0.51${\pm}$0.06 & 0.83${\pm}$0.03 \tabularnewline
\raggedright Random Forest          & 0.83${\pm}$0.04 & 0.64${\pm}$0.05 & 0.70${\pm}$0.04 & 0.51${\pm}$0.06 & 0.87${\pm}$0.03 \tabularnewline
\raggedright Decision Tree          & 0.76${\pm}$0.05 & 0.66${\pm}$0.06 & 0.68${\pm}$0.05 & 0.44${\pm}$0.09 & 0.71${\pm}$0.04 \tabularnewline
\raggedright Logistic regression    & 0.74${\pm}$0.04 & 0.69${\pm}$0.06 & 0.70${\pm}$0.04 & 0.44${\pm}$0.08 & 0.82${\pm}$0.03 \tabularnewline
\raggedright SVM                    & 0.73${\pm}$0.05 & 0.70${\pm}$0.05 & 0.70${\pm}$0.04 & 0.43${\pm}$0.08 & 0.81${\pm}$0.03 \tabularnewline
\raggedright Naive Bayes            & 0.74${\pm}$0.04 & 0.69${\pm}$0.05 & 0.69${\pm}$0.04 & 0.43${\pm}$0.07 & 0.79${\pm}$0.04 \tabularnewline
\raggedright KNN                    & 0.71${\pm}$0.05 & 0.67${\pm}$0.06 & 0.67${\pm}$0.04 & 0.38${\pm}$0.07 & 0.74${\pm}$0.04 \tabularnewline
\hline
\multicolumn{6}{@{}l}{\emph{Baselines}} \tabularnewline
\raggedright Vocab.\ threshold      & 0.66${\pm}$0.04 & 0.87${\pm}$0.03 & 0.74${\pm}$0.02 & 0.40${\pm}$0.07 & 0.77${\pm}$0.03 \tabularnewline
\raggedright Structure only         & 0.68${\pm}$0.04 & 0.67${\pm}$0.06 & 0.66${\pm}$0.04 & 0.36${\pm}$0.07 & 0.73${\pm}$0.04 \tabularnewline
\raggedright Random                 & 0.52${\pm}$0.05 & 0.51${\pm}$0.05 & 0.50${\pm}$0.04 & 0.02${\pm}$0.08 & 0.53${\pm}$0.05 \tabularnewline
\raggedright Majority               & 0.50${\pm}$0.00 & 1.00${\pm}$0.00 & 0.67${\pm}$0.00 & 0.00${\pm}$0.00 & 0.50${\pm}$0.00 \tabularnewline
\hline\hline
\end{tabular}
\end{table}

Error profiles differ across models. The tree ensembles are the most precise (Extra Trees: precision 0.87 and recall 0.63), that is, when they flag a test as flaky they are usually right, but they miss some of them. The vocabulary threshold rule does the opposite: it has the highest recall (0.87) but a precision of only 0.66, lower than that of every trained classifier, because it flags as flaky any test containing the most informative tokens. The linear models, Naive Bayes, and KNN improve little or not at all over the structure-only model. Since the folds are balanced, these precisions do not transfer to real use, which the cross-repository validation quantifies.
In cross-repository validation (Table 4), which uses the natural flaky rate (0.39\%), AUC drops to 0.75--0.78 and precision is only 2--3\% for the two best models: for every flaky test found, there are dozens of false alarms. The average precision (AP) of the two best models is four to six times the natural rate, which serves as the chance reference. This shows that the ranking is informative, but that automatic gating would be unusable. On this metric, the vocabulary rule is competitive (0.025, against 0.015 and 0.023 for the two models), although the tree ensembles lead on AUC.

\begin{table}[!h]
\caption{Cross-repository validation (LOPO) at the natural rate (0.39\%). Macro: mean over the nine repositories; pooled: all held-out tests together.}
\label{tab:lopo}
\centering
\scriptsize
\setlength{\tabcolsep}{3.5pt}
\begin{tabular}{@{}lrrrrrr@{}}
\hline\hline
 & \multicolumn{4}{c}{\textbf{Macro}} & \multicolumn{2}{c}{\textbf{Pooled}} \\
\textbf{Model} & \textbf{AUC} & \textbf{AP} & \textbf{Rec.} & \textbf{Prec.} & \textbf{AUC} & \textbf{AP} \\
\hline
Extra Trees          & 0.757 & 0.055 & 0.376 & 0.028 & 0.775 & 0.015 \\
Random Forest        & 0.754 & 0.089 & 0.376 & 0.021 & 0.753 & 0.023 \\
Gradient Boosting    & 0.709 & 0.048 & 0.509 & 0.012 & 0.734 & 0.021 \\
Logistic regression  & 0.679 & 0.065 & 0.574 & 0.009 & 0.741 & 0.017 \\
SVM                  & 0.660 & 0.048 & 0.556 & 0.010 & 0.730 & 0.015 \\
Naive Bayes          & 0.644 & 0.053 & 0.567 & 0.009 & 0.718 & 0.018 \\
KNN                  & 0.597 & 0.020 & 0.529 & 0.008 & 0.711 & 0.011 \\
Decision Tree        & 0.597 & 0.011 & 0.476 & 0.012 & 0.740 & 0.010 \\
\hline
\multicolumn{7}{@{}l}{\emph{Baselines}} \\
Vocabulary threshold  & 0.682 & 0.040 & 0.732 & 0.008 & 0.721 & 0.025 \\
Structure only        & 0.674 & 0.020 & 0.559 & 0.009 & 0.654 & 0.006 \\
Random                & 0.516 & 0.013 & 0.507 & 0.005 & 0.554 & 0.006 \\
Majority              & 0.500 & 0.005 & 1.000 & 0.005 & 0.500 & 0.004 \\
\hline\hline
\end{tabular}
\end{table}

Table 5 shows the variation across repositories for the two best models. AUC ranges from 0.53 (NetMQ, with 2 flaky tests) to 0.97 (Testcontainers, with 1); repositories with a single flaky test (Akka.NET, Quartz.NET, and Testcontainers) are only illustrative. Among those with several flaky tests, Nerdbank.Streams is the most favorable (AUC of 0.74 to 0.80 and AP of 0.15 to 0.18), while MassTransit and StreamJsonRpc, which together hold 14 of the 30 flaky tests, have AUCs of 0.67 to 0.70. Figure 3 summarizes the main effect: for every model, AUC drops when the test's repository is not part of training; the drop is smaller for the tree ensembles (about 0.10 to 0.12) than for logistic regression, SVM, Naive Bayes, and KNN (0.14 to 0.15).

\begin{table}[!h]
\caption{Per-repository LOPO for the two best models (RF: Random Forest; ET: Extra Trees). Repositories with a single flaky test are only illustrative.}
\label{tab:lopo-repo}
\centering
\scriptsize
\setlength{\tabcolsep}{4pt}
\begin{tabular}{@{}lrrrrrr@{}}
\hline\hline
 & & & \multicolumn{2}{c}{\textbf{RF}} & \multicolumn{2}{c}{\textbf{ET}} \\
\textbf{Repository} & \textbf{Flaky} & \textbf{Tests} & \textbf{AUC} & \textbf{AP} & \textbf{AUC} & \textbf{AP} \\
\hline
Akka.NET         & 1 & 1,134 & 0.782 & 0.006 & 0.862 & 0.008 \\
MassTransit      & 6 & 1,548 & 0.672 & 0.023 & 0.704 & 0.019 \\
Nerdbank.Streams & 6 &   595 & 0.795 & 0.179 & 0.741 & 0.154 \\
NetMQ            & 2 &   259 & 0.530 & 0.015 & 0.549 & 0.017 \\
NLog             & 2 & 1,692 & 0.844 & 0.009 & 0.834 & 0.009 \\
Polly            & 3 &   559 & 0.775 & 0.033 & 0.846 & 0.041 \\
Quartz.NET       & 1 & 1,016 & 0.737 & 0.004 & 0.646 & 0.003 \\
StreamJsonRpc    & 8 &   763 & 0.681 & 0.020 & 0.696 & 0.022 \\
Testcontainers   & 1 &   192 & 0.969 & 0.507 & 0.933 & 0.223 \\
\hline\hline
\end{tabular}
\end{table}

\begin{figure}[!h]
\centering
\footnotesize
\setlength{\unitlength}{1pt}
\begin{picture}(248,138)(0,-14)
\thinlines
\put(21,0.0){\line(1,0){3}}
\put(19,0.0){\makebox(0,0)[r]{0.5}}
\put(21,27.5){\line(1,0){3}}
\put(19,27.5){\makebox(0,0)[r]{0.6}}
\put(21,55.0){\line(1,0){3}}
\put(19,55.0){\makebox(0,0)[r]{0.7}}
\put(21,82.5){\line(1,0){3}}
\put(19,82.5){\makebox(0,0)[r]{0.8}}
\put(21,110.0){\line(1,0){3}}
\put(19,110.0){\makebox(0,0)[r]{0.9}}
\put(24,0){\line(1,0){220}}
\put(24,0){\line(0,1){116}}
\put(24,119){\makebox(0,0)[b]{AUC}}
\put(29.8,0){\rule{8pt}{99.0pt}}
\put(37.8,0){\framebox(8,70.7){}}
\put(37.8,-8){\makebox(0,0){ET}}
\put(57.2,0){\rule{8pt}{101.8pt}}
\put(65.2,0){\framebox(8,69.8){}}
\put(65.2,-8){\makebox(0,0){RF}}
\put(84.8,0){\rule{8pt}{90.7pt}}
\put(92.8,0){\framebox(8,57.5){}}
\put(92.8,-8){\makebox(0,0){GB}}
\put(112.2,0){\rule{8pt}{88.0pt}}
\put(120.2,0){\framebox(8,49.2){}}
\put(120.2,-8){\makebox(0,0){LR}}
\put(139.8,0){\rule{8pt}{85.3pt}}
\put(147.8,0){\framebox(8,44.0){}}
\put(147.8,-8){\makebox(0,0){SVM}}
\put(167.2,0){\rule{8pt}{79.8pt}}
\put(175.2,0){\framebox(8,39.6){}}
\put(175.2,-8){\makebox(0,0){NB}}
\put(194.8,0){\rule{8pt}{66.0pt}}
\put(202.8,0){\framebox(8,26.7){}}
\put(202.8,-8){\makebox(0,0){KNN}}
\put(222.2,0){\rule{8pt}{57.7pt}}
\put(230.2,0){\framebox(8,26.7){}}
\put(230.2,-8){\makebox(0,0){DT}}
\put(40,116){\rule{6pt}{6pt}}
\put(49,119){\makebox(0,0)[l]{Cross-validation}}
\put(130,116){\framebox(6,6){}}
\put(139,119){\makebox(0,0)[l]{Cross-repository (macro)}}
\end{picture}
\caption{AUC under cross-validation (balanced folds) and macro cross-repository AUC, by model (the axis starts at 0.5, the chance level; ET: Extra Trees; RF: Random Forest; GB: Gradient Boosting; LR: logistic regression; NB: Naive Bayes; DT: decision tree).}
\label{fig:auc}
\end{figure}

\noindent\fbox{\parbox{\dimexpr\columnwidth-2\fboxsep-2\fboxrule\relax}{\textbf{RQ1.1 and RQ2:} vocabulary predicts flakiness in C\# well above chance (MCC 0.51--0.53, AUC 0.86--0.87 under cross-validation; AUC 0.75--0.78 across repositories), but the gain over structure is about $+$0.15 MCC and only for the tree ensembles. The practical value lies in ranking.}}

\subsection{RQ1.2: label robustness}
The two rerun rounds diverge: of the 36 flaky units found by either of them, only 13 are flaky in both (Jaccard 0.36); 17 are flaky only in the 100-run round and 6 only in the 20-run round. The 20-run round finds 43\% of what the 100-run round finds, and the 100-run round finds 68\% of what the 20-run round finds. Most tests exclusive to the 100-run round fail rarely (for example, Nerdbank.Streams 6$\to$0, NetMQ 2$\to$0), and of the six exclusive to the 20-run round, four are in NLog and three of them failed together in a single run, so they represent a single event. One of the tests exclusive to the 100-run round failed in all 20 manual runs, which suggests an environmental difference between the rounds. In LOPO with the 20-run labels, only five repositories have flaky tests, and the best macro AUC is 0.72. With the 20-run labels (Table 7), the tree ensembles lose 0.07 to 0.14 MCC and 0.04 to 0.12 AUC, and the structure-only model (MCC 0.47) catches up with them within the confidence interval; the advantage of vocabulary holds only with the richer labels.

\begin{table}[!h]
\caption{Overlap between the flaky units of the two rounds.}
\label{tab:sobreposicao}
\centering
\footnotesize
\begin{tabular}{@{}lcp{0.50\columnwidth}@{}}
\hline\hline
\textbf{Category} & \textbf{Units} & \textbf{Where} \tabularnewline
\hline
Flaky in both & 13 & \raggedright StreamJsonRpc 5, MassTransit.Tests 3, Polly 2, Akka.NET 1, MassTransit.Abstractions 1, NLog 1 \tabularnewline[2pt]
Only in 100 runs & 17 & \raggedright Nerdbank.Streams 6, StreamJsonRpc 3, NetMQ 2, MassTransit.Tests 2, Testcontainers 1, Quartz 1, NLog 1, Polly 1 \tabularnewline[2pt]
Only in 20 runs & 6 & \raggedright NLog 4, MassTransit.Tests 1, Polly 1 \tabularnewline
\hline\hline
\end{tabular}
\end{table}

\begin{table}[!h]
\caption{Cross-validation with 100-run vs.\ 20-run labels (19 flaky units in the latter).}
\label{tab:cv100x20}
\centering
\footnotesize
\setlength{\tabcolsep}{5pt}
\begin{tabular}{@{}lcccc@{}}
\hline\hline
 & \multicolumn{2}{c}{\textbf{MCC}} & \multicolumn{2}{c}{\textbf{AUC}} \\
\textbf{Model} & \textbf{100$\times$} & \textbf{20$\times$} & \textbf{100$\times$} & \textbf{20$\times$} \\
\hline
Extra Trees       & 0.53 & 0.43 & 0.86 & 0.82 \\
Random Forest     & 0.51 & 0.44 & 0.87 & 0.80 \\
Gradient Boosting & 0.51 & 0.37 & 0.83 & 0.71 \\
\hline
\multicolumn{5}{@{}l}{\emph{Baselines}} \\
Vocab.\ threshold & 0.40 & 0.23 & 0.77 & 0.66 \\
Structure only    & 0.36 & 0.47 & 0.73 & 0.77 \\
Random            & 0.02 & 0.04 & 0.53 & 0.54 \\
\hline\hline
\end{tabular}
\end{table}

As an independent check, the six tests flaky only in the 20-run round are labeled stable under the 100-run labels and receive a score when their own repository is held out, so the model has never seen them as positives. Three of them rank in the top 15\% of their repository (percentiles 0.96, 0.95, and 0.85), and the three NLog tests that failed together fall at percentiles 0.30--0.42. Counting the trio as a single event, three of four independent events are ranked high.

\medskip
\noindent\fbox{\parbox{\dimexpr\columnwidth-2\fboxsep-2\fboxrule\relax}{\textbf{RQ1.2:} the signal is robust in direction but not in magnitude: the round that produced the labels changes which tests are flaky (Jaccard 0.36), and the relative value of vocabulary over structure depends on it.}}
\medskip

\subsection{RQ3: the vocabulary of flaky tests in C\#}
Table 8 shows the tokens with the highest information gain that occur in flaky tests from at least three repositories. The most informative markers are Task-based concurrency, that is, operations run in the background and awaited with async/await (task in 80\% of flaky tests vs.\ 31\% of stable ones; async/await in 67\% vs.\ 30\%), time handling (TimeSpan, CancellationToken, Sleep, Delay), and resource lifecycles (Dispose, StopAsync), along with rare, specific markers such as GC and WeakReference (Table 9). This agrees with Java, JavaScript, and Swift and with the sleep/timeout pattern described for .NET \cite{leesatapornwongsa2022flakerepro}. The frequent markers are general: task, async, and await occur in six repositories. However, a third of stable tests in C\# also contain async/await, which explains why presence rules are weak.

\begin{table}[!h]
\caption{Top tokens by information gain (IG, in nats) present in flaky tests from at least three repositories. \%F and \%S: proportion of flaky (30) and stable (12,635) units containing the token.}
\label{tab:tokens}
\centering
\footnotesize
\setlength{\tabcolsep}{12pt}
\begin{tabular}{@{}lrr@{}}
\hline\hline
\textbf{Token} & \textbf{IG} & \textbf{\%F} \\
\hline
\texttt{task}              & 0.145 & 80 \\
\texttt{async task}        & 0.084 & 67 \\
\texttt{await}             & 0.084 & 67 \\
\texttt{async}             & 0.082 & 67 \\
\texttt{length}            & 0.065 & 27 \\
\texttt{await task}        & 0.061 & 23 \\
\texttt{server}            & 0.056 & 23 \\
\texttt{enumerable}        & 0.044 & 13 \\
\texttt{logger}            & 0.044 & 20 \\
\texttt{fromseconds}       & 0.043 & 20 \\
\texttt{writeline}         & 0.038 & 17 \\
\texttt{timespan}          & 0.034 & 23 \\
\texttt{dispose}           & 0.034 & 13 \\
\texttt{cancellationtoken} & 0.025 & 13 \\
\hline\hline
\end{tabular}
\end{table}

The raw information-gain ranking is dominated by project conventions: trait, fact trait, and checkgcpressureasync come only from StreamJsonRpc, and trait originates from an attribute, not from the test body. Information gain favors frequent tokens; we therefore also report rare but specific markers by lift (Table 9). They fall into memory and garbage collection (gc, weakreference), asynchronous composition (whenall), cancellation (cancellationtokensource), lifecycles (stopasync, disposed), and fixed waits (sleep, task delay). Each occurs in only two to four flaky tests, so they are exploratory.

\begin{table}[!h]
\caption{Rare, specific markers (at least two flaky tests in at least two repositories, except \texttt{weakreference}). $p$-values: one-sided Fisher's exact tests, uncorrected for the roughly 48,000 candidate tokens.}
\label{tab:marcadores}
\centering
\scriptsize
\setlength{\tabcolsep}{5pt}
\begin{tabular}{@{}lrrrrr@{}}
\hline\hline
\textbf{Token} & \textbf{Flaky} & \textbf{\%F} & \textbf{\%S} & \textbf{\emph{Lift}} & \textbf{\emph{p}} \\
\hline
\texttt{gc}                      & 3 & 10.0 & 0.12 & 84$\times$ & $<$0.001 \\
\texttt{weakreference}           & 2 &  6.7 & 0.18 & 37$\times$ & 0.002 \\
\texttt{stopasync}               & 2 &  6.7 & 0.12 & 56$\times$ & --- \\
\texttt{disposed}                & 2 &  6.7 & 0.13 & 50$\times$ & --- \\
\texttt{whenall}                 & 4 & 13.3 & 0.62 & 22$\times$ & --- \\
\texttt{cancellationtokensource} & 4 & 13.3 & 1.56 &  9$\times$ & --- \\
\texttt{sleep}                   & 2 &  6.7 & 0.82 &  8$\times$ & 0.03 \\
\texttt{task delay}              & 2 &  6.7 & 0.97 &  7$\times$ & 0.04 \\
\hline\hline
\end{tabular}
\end{table}

Relation to root causes. The markers connect to the causes of flakiness described in the Introduction. task, async, and await indicate background operations whose completion the test must await, the territory of race conditions. TimeSpan, FromSeconds, CancellationToken, Sleep, and Delay indicate dependence on time and timeouts. Dispose, StopAsync, server, logger, and WriteLine indicate external resources that must be created and released, such as servers and output streams. GC and WeakReference indicate dependence on the garbage collector, whose timing is nondeterministic. This reading is an interpretation of the tokens, not a confirmed classification of each test's cause.
Stability markers. Unlike Swift, where xctassertequal marks clearly synchronous tests, we found no robust stability marker: among the tokens present in flaky tests from at least three repositories, the only one more frequent in stable tests is assert (67\% vs.\ 69\%), which does not discriminate. The token areequal (0\% vs.\ 20\%) is an NUnit assertion style used by repositories without flaky tests, so it reflects the choice of framework, not stability.

\medskip
\noindent\fbox{\parbox{\dimexpr\columnwidth-2\fboxsep-2\fboxrule\relax}{\textbf{RQ3:} flakiness in C\# is signaled by \texttt{Task}-based concurrency, time handling, and resource lifecycles, complemented by rare markers such as \texttt{GC} and \texttt{WeakReference}. The frequent markers are general across repositories, the rare ones are specific but backed by few tests, and no robust stability marker emerged.}}
\medskip

\subsection{RQ4: error analysis}
To analyze errors, we use the LOPO scores, in which each test is scored by models trained without its repository. The score is the mean of Random Forest and Extra Trees over the 10 undersampling draws. In each of the nine repositories that have flaky tests, we rank all labeled tests (flaky and stable) by score, breaking ties by average rank, and compute each test's percentile within its repository, so that 1.0 corresponds to the highest-scoring test. The 30 flaky tests (Table 10) have a median percentile of 0.80; 9 have a percentile of at least 0.90, that is, they are in the top 10\% of their repository, and 23 have a percentile of at least 0.75, in the top 25\%. This means that, if each repository were ranked by the model and the top quartile were selected for rerunning (about a quarter of the tests), that quartile would contain 77\% of the flaky tests, against about 25\% expected by chance, with a 95\% confidence interval of 58\% to 90\%. The result also does not depend on the repositories that concentrate the positives: in the three largest (20 of the 30 flaky tests), 14 are in the top quartile (70\%), and in the other six repositories, 9 of 10 (90\%), although these have only one or two each.

\begin{table}[!h]
\caption{Position of the 30 flaky tests in their repositories' rankings (LOPO; percentile 1.0 = highest score).}
\label{tab:ranking}
\centering
\footnotesize
\begin{tabular}{@{}p{0.48\columnwidth}cl@{}}
\hline\hline
\textbf{Criterion} & \textbf{Flaky} & \textbf{Proportion} \tabularnewline
\hline
\raggedright Percentile $\geq$ 0.90 (top 10\%) & 9 & 30\% \tabularnewline
\raggedright Percentile $\geq$ 0.75 (top 25\%) & 23 & 77\% (CI 58--90\%) \tabularnewline
\raggedright Percentile $<$ 0.50 & 6 & 20\% \tabularnewline
\raggedright Three largest repositories, top quartile & 14 of 20 & 70\% \tabularnewline
\raggedright Other six repositories, top quartile & 9 of 10 & 90\% \tabularnewline
\raggedright Median percentile & 0.80 & \tabularnewline
\hline\hline
\end{tabular}
\end{table}

This estimate should be read with caution. The two models that make up the score were chosen for their performance on the same data, and the confidence interval treats tests as independent, even though they come from a convenience sample and some are correlated (three NLog tests failed together), so the real uncertainty is probably larger. We also assume the same execution cost per test, which probably underestimates the real cost, since asynchronous tests and tests with fixed waits tend to be slower. Moreover, the result measures how concentrated flaky tests are at the top of the ranking, not the chance that rerunning reveals the flakiness, which depends on how often each test fails.
We also examined, in the test source code, the errors for the tests the model did not place in the top quartile (seven of the 30) and for the highest-scoring stable tests. This is a qualitative analysis of a few examples, not a count of causes. The false negatives are flaky tests whose cause lies beyond the reach of vocabulary. Three patterns recur. (i) The test merely delegates to a helper method: JsonRpcPerf\_Pipe and JsonRpcPerf\_Channel (Nerdbank.Streams) are one-line calls to a shared helper, with no assertions, so the cause lies outside the test body; InvokeWithProgressParameter\_NoMemoryLeakConfirm (StreamJsonRpc) awaits a helper, forces GC.Collect(), and checks a WeakReference, and its nondeterminism comes from garbage collection. (ii) Fixed waits without asynchronous constructs: ReconnectOnRouterBug (Section 2) synchronizes only through Thread.Sleep, a token present in only two flaky tests and weakly learned. (iii) Time and global state: MassTransit's Should\_be\_using\_the\_correct\_algorithm\_for\_sequential\_guids installs a tick provider and compares parts of consecutively generated identifiers, which depends on the clock and, apparently, on shared static state.
The false positives are stable tests that use concurrency and time deterministically. Polly's Execute\_CallerCancellation\_EnsureExceptionCarriesCallerToken uses an injected time provider and a cancellation token, being asynchronous and time-related, yet deterministic by construction; Quartz's Trigger\_Fires\_After\_Clock\_Jumps\_Backward\_And\_Returns controls time through a custom provider and awaits a TaskCompletionSource instead of sleeping. The model detects the constructs but does not judge whether they are used well. Akka.NET's A\_CircuitBreaker\_stress\_test uses a fixed await Task.Delay(3000) with a wide margin and may be latently flaky beyond our 120 runs, so some false positives may be label noise.

\medskip
\noindent\fbox{\parbox{\dimexpr\columnwidth-2\fboxsep-2\fboxrule\relax}{\textbf{RQ4:} the model fails when the cause is hidden in helper methods, in GC behavior, or in shared state (false negatives), and when asynchronous or time constructs are used with controlled time or event-based synchronization (false positives). Judging whether usage is correct requires semantic, not lexical, analysis.}}
\medskip

\subsection{Sensitivity analyses}
We tested two threats to the vocabulary signal, keeping the labels, undersampling seeds, and folds fixed so that the comparison is paired. Without attributes: we removed all tokens inside attribute lists ([Fact], [Trait(...)], [InlineData(...)]), so that project conventions do not drive the prediction. With helpers: we added to each test the tokens of the helper methods it calls, resolved by name (at most five helpers per test, one level only); the expansion resolved at least one helper in 67\% of the methods, and, since resolution is by name, without type information, some helpers may merely share a name with the one actually called.

\begin{table}[!h]
\caption{Sensitivity of cross-validation MCC and LOPO macro AUC (original features, without attributes, with one-level helpers).}
\label{tab:sensibilidade}
\centering
\scriptsize
\setlength{\tabcolsep}{3.5pt}
\begin{tabular}{@{}lcccccc@{}}
\hline\hline
 & \multicolumn{3}{c}{\textbf{MCC (CV)}} & \multicolumn{3}{c}{\textbf{Macro AUC (LOPO)}} \\
\textbf{Model} & \textbf{orig.} & \textbf{no attr.} & \textbf{+help.} & \textbf{orig.} & \textbf{no attr.} & \textbf{+help.} \\
\hline
Extra Trees         & 0.53 & 0.46 & 0.56 & 0.757 & 0.737 & 0.764 \\
Random Forest       & 0.51 & 0.46 & 0.53 & 0.754 & 0.721 & 0.773 \\
Gradient Boosting   & 0.51 & 0.49 & 0.43 & 0.709 & 0.683 & 0.643 \\
Logistic regression & 0.44 & 0.43 & 0.49 & 0.679 & 0.673 & 0.714 \\
\hline
\multicolumn{7}{@{}l}{\emph{Baseline}} \\
Structure only      & 0.36 & 0.36 & 0.36 & 0.674 & 0.674 & 0.674 \\
\hline\hline
\end{tabular}
\end{table}

Removing attributes reduces the MCC of the two best ensembles from 0.53 to 0.46 (Extra Trees) and from 0.51 to 0.46 (Random Forest), but keeps AUC at 0.84--0.85: attributes carry part of the signal, plausibly project conventions, but most of it lies in the test body, and the ensembles remain well above the structure-only baseline. Adding helpers slightly improves the ensembles (MCC 0.56 and 0.53) and improves the linear and probabilistic models more, but hurts Gradient Boosting (from 0.51 to 0.43) and the Decision Tree (from 0.44 to 0.39). The ensembles' gains are within the confidence intervals. The naive one-level expansion therefore does not eliminate the helper blind spot: it adds tokens to two thirds of the tests, flaky or stable alike.

\section{Discussion}
Table 12 summarizes the answers to the research questions; the following subsections discuss what they mean.

\begin{table}[!h]
\caption{Summary of answers to the research questions.}
\label{tab:resumo}
\centering
\footnotesize
\begin{tabular}{@{}lp{0.36\columnwidth}p{\dimexpr0.64\columnwidth-4\tabcolsep-2em\relax}@{}}
\hline\hline
\textbf{RQ} & \textbf{Answer} & \textbf{Main evidence} \tabularnewline
\hline
1.1 & \raggedright Above chance, modestly & \raggedright MCC 0.51--0.53 and AUC 0.86--0.87 (CV); AUC 0.75--0.78 and precision 2--3\% (LOPO) \tabularnewline[2pt]
1.2 & \raggedright Stable in direction, not in magnitude & \raggedright 13 of 36 flaky units in both rounds; \emph{ensembles} lose 0.07--0.14 MCC \tabularnewline[2pt]
2 & \raggedright Moderate margin over structure and threshold & \raggedright Structure only: MCC 0.36; threshold: 0.40; vocabulary: $+$0.15--0.17 for the best models \tabularnewline[2pt]
3 & \raggedright \texttt{Task}, time, and lifecycles; no stability marker & \raggedright \texttt{task} in 80\% of flaky vs.\ 31\% of stable tests; \texttt{GC} 84$\times$ and \texttt{WeakReference} 37$\times$ (rare) \tabularnewline[2pt]
4 & \raggedright Fails when the cause lies outside the body or usage is deterministic & \raggedright 23 of 30 in the top quartile; 7 outside it \tabularnewline
\hline\hline
\end{tabular}
\end{table}

Rank, don't gate. What the model does well is order tests, not decide whether each one is flaky. When it flags a test as flaky, it is right only 2 or 3 times out of 100, so using it to block a delivery in the pipeline would raise too many false alarms. Its usefulness appears when we consider a team that lacks the time to rerun every test and must choose where to start. In that case, the ranking works as a priority list: the top quarter of each project's list concentrates 77\% of the flaky tests (Section 6.5), and rerunning that quarter would cost about a quarter of the total effort, if each test took roughly the same time. An example helps show the gain: in a repository of 1,000 tests at the natural rate of 0.39\%, there are about 4 flaky tests; the top quartile has 250 tests and, with 77\% coverage, about 3 of them, a density of roughly 1.2\%, three times that of the whole set, though still low for gating. This holds as long as rerunning actually reveals the flakiness, which depends on how often each test fails. There is also a secondary benefit: the vocabulary analysis shows which constructs tend to accompany flakiness, such as fixed waits (Sleep), cancellation, and garbage collection, and a code reviewer can use these cues as warnings during review. We did not test this use with practitioners. The ranking can also select which tests to hand to dynamic tools that reproduce concurrency failures, such as FlakeRepro [13], which need the test and an observed failure to work.
Part of the signal is project-specific. When the model is trained and tested on tests from all projects mixed together, it performs very well (AUC of 0.86--0.87). When it must predict the tests of a project it has never seen, performance drops (AUC of 0.75--0.78). This drop indicates that part of what the model learns is each project's writing habits, not flakiness itself. Three repositories hold two thirds of the flaky tests (20 of 30), removing attributes reduces MCC by about 0.05 to 0.07, and the tokens with the highest information gain come largely from a single project. On the other hand, cross-repository performance remains well above chance, which shows that a general signal also exists. In practice, this calls for caution when importing results from other teams and languages, and, if the tool is adopted in an organization, it is worth calibrating it with the flaky tests that organization already knows about.
Comparison with other ecosystems. In Swift, the same kind of classifier reached an MCC of 0.75 under balanced cross-validation, against 0.08 for the best baseline [15]; in C\#, the MCC is 0.51--0.53, and the non-trivial baselines reach 0.36--0.40. The numbers are not directly comparable, because the datasets, labels, and number of positives differ (91 flaky in Swift, 30 here), but they suggest that the flakiness vocabulary in C\# is less distinctive: asynchronous constructs are also common in stable tests, and there is no stability marker like Swift's xctassertequal.
Toward better features. The errors point to natural extensions: propagating tokens through the helper methods called (interprocedural analysis) and designing features that separate fixed waits from controlled time. A simple version of the first, which adds to each test the tokens of the helpers it calls, brought small and inconsistent gains (Section 6.6); a proper analysis, with type information, is left for future work. The second requires distinguishing, in the code, a Sleep that assumes a duration from a Delay controlled by a simulated clock, which vocabulary alone cannot do.
Methodological lessons. First, the unit of analysis changes the answer. A single test method can generate several cases that the runner reports separately: counting cases, there are 50 flaky tests; counting the methods that declare them, there are 30. In StreamJsonRpc, for example, the same test runs with six serializers and generates 25 flaky cases from only 8 methods. Second, labels depend on the round of runs: of the 36 flaky units found by either of the two rounds, only 13 are flaky in both, so studies with a single round should report this sensitivity. Third, we document .NET-specific pitfalls that distorted our initial counts (Table 1): parameterized test names containing timestamps, thread ids, or dates, which make the same test look like several; differences between target frameworks; tests declared in base classes and executed in several derived classes; and contaminated environments, which produced spurious flaky tests.
Recommendations for rerun studies in .NET. The pitfalls we encountered suggest a minimal protocol for anyone wishing to repeat this kind of study: (i) identify each test by its fully qualified class plus its display name and normalize timestamps, thread ids, and dates in the names; (ii) resolve inherited and parameterized tests to the declaring method and report both counts, per case and per method; (iii) build once, then run without rebuilding and on the same target framework in every round; (iv) isolate the environment between runs and discard contaminated rounds; and (v) report the sensitivity of the labels to the number of runs.

\section{Threats to Validity}
Labels. The labels come only from reruns, so tests whose flakiness does not reproduce on our machine are labeled stable, and the number of runs matters: the 20-run round identified only 13 of the 30 flaky tests identified with 100 runs. Environmental disruptions can also create spurious flaky tests; we mitigated them by discarding contaminated rounds and repeating them in isolation. Test identity is reconstructed from display names and class hierarchies, and a small number of tests could not be resolved to source code (none of them flaky in the 100-run round). Internal. The 1:1 undersampling balances the test folds and makes precision and F1 optimistic (LOPO quantifies the difference); the same 30 positives appear in every draw, so the confidence intervals underestimate uncertainty and, with so few positives, the ordering among classifiers is not reliable; correlated tests (for example, three NLog tests that failed together) may inflate the positive class; and the two best models were chosen on the same data on which they are evaluated. Attribute tokens may encode project conventions, and removing them reduces the MCC of the best models by about 0.05-0.07. External and selection. Project selection was not random: it favored open-source libraries with unit tests and domains in which concurrency is common, so the 0.24\% rate and the results may not generalize to applications, UI tests, or proprietary code; moreover, flakiness is concentrated in nine repositories, and the results are specific to the .NET toolchain and machine used.

\section{Conclusion}
This work evaluated static prediction of flaky tests in C\# from code vocabulary. The lexical signal exists, but it is modest: across the 17 projects, which total 19,211 executed test cases, corresponding to 12,665 labeled test methods, of which 30 are flaky, the best models reach an MCC of 0.51--0.53 and an AUC of 0.86--0.87 under cross-validation, and an AUC of 0.75--0.78 when the repository is unseen, with a precision of 2--3\% at the natural base rate. The signal rests on Task-based concurrency, that is, operations run in the background and awaited with async/await, on time handling, and on resource lifecycles, and it is useful for prioritizing reruns, not for gating pipelines. Seven of the 30 flaky tests fall outside the top quartile of their repository's ranking, meaning the model does not place them among those that would be rerun first. Reading the code of some of these tests, we noticed that, in certain cases, the cause of flakiness is not in the test body, which is what the vocabulary sees: the test merely calls a helper method where the wait happens, or depends on garbage collector behavior or on state shared between tests. The model also gives high scores to stable tests that use time and concurrency in a controlled way, such as waiting for events and simulated clocks, because the words they are made of resemble those of flaky tests, even though their usage is correct. This is a qualitative analysis of a few examples, not a count of causes. As next steps, we plan to make the model follow calls between methods, so that it also reads the code of the helper methods the test uses; to design features that distinguish a fixed wait from time controlled by the test itself; to repeat the study with more projects, chosen at random; and to ask practitioners how they would use a test ranking to decide what to rerun.

% Required by the MSR 2027 Open Science Policy.
\section*{Data Availability}
The scripts used to execute the test suites and collect the reruns, the
resulting dataset of flaky and stable test methods, the feature extraction
and classification pipeline, and instructions to reproduce all results are
available in an anonymized replication package:
\url{https://anonymous.4open.science/r/SUBSTITUIR-PELO-LINK-REAL}.
Upon acceptance, the package will be archived on Zenodo under the CC-BY~4.0
license (data) and an open-source license (code).

\bibliography{biblio}
\bibliographystyle{IEEEtran}

\end{document}